\section{Mode Collapse and Verified Replay}
\label{app:theory}

\noindent\emph{Main-body link.} This appendix supplies the formal claims
invoked in Sections~\ref{sec:method} and~\ref{sec:results}.

A binary verifier assigns the same reward to every correct mode, but an
update can still amplify the most common one. We first explain that mechanism
in a finite categorical model, then ask what changes when we alter the update
geometry, add entropy, or replay earlier successes. Replay directly protects
represented modes; discovery, admission, and reliable neural optimization
remain separate requirements.

\paragraph{How to read the proof chain.}
Begin with the three-correct-mode example in
Section~\ref{app:theory-grpo-mean}. The group calculations reveal why a
correctness gradient need not balance correct alternatives, and the next
section turns the local imbalance into an asymptotic collapse theorem.
Sections~\ref{app:theory-size}--\ref{app:theory-entropy-comparison} then vary
geometry and entropy: they show both how collapse can slow and why it is not
a universal consequence of rewarding correctness.

The second part asks what replay adds. Sections~\ref{app:theory-gradient-availability}
and~\ref{app:theory-replay} separate a missing fresh signal from a persistent
bank signal, then turn bounded replay loss into protected probabilities.
The remaining sections connect that argument to complete responses,
controlled stochastic updates, and actual bank admission. Finally,
Section~\ref{app:theory-survival-certificates} shows how much those bounds
can certify and what finite observations can establish about a named mode.
Each proof states its assumptions and attributes the literature tools it uses.

All worked numbers below are illustrative calculations in the stated models,
not measurements from the training runs. Unless a new example is introduced,
the examples reuse the same three correct modes and one incorrect category.
Comparisons keep the optimized objective, decoding law, and time variable
explicit, since changing any of these can change the conclusion.

Here $G$ counts training rollouts, $K$ counts evaluation draws, $m,n$ count
correct and incorrect categories, and $k$ is the bank size. In the mean-flow results, $t$ is
continuous optimization time and $\tau$ is a reparameterized learning clock.
The later stochastic and admission results use explicitly defined update or
opportunity indices. None is a wall-clock or GPU-runtime prediction. ``Collapse'' refers to
limiting concentration: finite softmax logits already give positive
probabilities at every finite time. A uniform lower bound for all future
times is the stronger retention property proved below.

